\section{Original NDL Architecture and Claims}

\subsection{Architecture and growth rule}
\begin{draftpoints}
  \item published model is a symmetric stacked autoencoder whose encoder widths decrease toward a latent bottleneck and whose decoder mirrors those widths \parencite{draelos2017neurogenesis}
  \item MNIST base training uses digits 1 and 7 with encoder widths $[200,100,75,20]$; the remaining digits arrive sequentially in the order $[0,2,3,4,5,6,8,9]$ \parencite{draelos2017neurogenesis}
  \item For each depth, samples whose global reconstruction error exceeds a stored threshold are treated as outliers; growth continues until the allowed outlier count is reached or a maximum growth allowance is exhausted
  \item During the plasticity phase, mature encoder weights are frozen, new encoder weights use the base learning rate, and decoder weights use one hundredth of that rate; a stability phase then trains all weights in the current local autoencoder at the reduced rate \parencite{draelos2017neurogenesis}
  \item New lower-level columns must be propagated into the next layer so the expanded representation can influence deeper features
\end{draftpoints}

\begin{figure}[htbp]
\centering
\includegraphics[width=0.82\linewidth]{original_architecture.png}
\caption{Original NDL illustration reproduced from Draelos et al. The figure motivates neuron insertion when a novel input is reconstructed poorly and shows that reconstruction can be evaluated at several autoencoder depths \parencite{draelos2017neurogenesis}.}
\label{fig:original-architecture}
\end{figure}

\subsection{Intrinsic replay}
\begin{draftpoints}
  \item Intrinsic replay stores a class-conditional mean and Cholesky factor of the top-level latent covariance rather than retaining old training images \parencite{draelos2017neurogenesis}
  \item Replay samples are drawn from the fitted latent Gaussian and decoded through the current autoencoder; these generated approximations are mixed with the incoming class during stabilization
  \item Intrinsic replay must be evaluated separately from neuron growth because replay can improve retention even when architecture remains fixed
\end{draftpoints}

\begin{figure}[htbp]
\centering
\includegraphics[width=0.56\linewidth]{original_intrinsic_replay.png}
\caption{Intrinsic replay mechanism from the original paper: class-conditional latent statistics are sampled and decoded to approximate previously observed classes \parencite{draelos2017neurogenesis}.}
\label{fig:original-ir}
\end{figure}

\subsection{Published experimental claims}
\begin{draftpoints}
  \item central MNIST claim is directional: NDL with intrinsic replay is reported to reconstruct learned digits slightly better than a fixed, capacity-matched conventional learner with intrinsic replay \parencite{draelos2017neurogenesis}
  \item paper also reports that growth contributes useful new capacity, preserves previously learned representations, and remains beneficial in the larger NIST SD-19 experiment
  \item publication does not fully specify optimizer family, activation functions, thresholds, outlier quota, phase-length interpretation, or the new-node-count policy, preventing binary-equivalent reconstruction from the text alone
\end{draftpoints}

\begin{table}[htbp]
\centering
\caption{Published claims and corresponding replication tests.}
\label{tab:claims-to-test}
\small
\begin{tabularx}{\textwidth}{p{3.2cm}YYp{2.5cm}}
\toprule
Claim & Original evidence & Replication measure & Decision rule \\
\midrule
NDL+IR exceeds CL+IR & Figure 4 reconstruction curves & Ten-seed macro MSE and forgetting & Direction and confidence intervals \\
Growth protects old classes & Curves and reconstruction images & Per-class MSE and positive forgetting & Lower old-class degradation \\
Capacity emerges from demand & Layer-growth plots & Quota stops, cap exhaustion, final widths & Demand satisfied before cap \\
Cost is manageable & Complexity discussion & Updates, runtime, CUDA and RSS peaks & Cost relative to matched CL \\
Effect extends to SD-19 & Figures 6--8 & Three-seed SD-19 MSE and forgetting & Same qualitative ordering \\
\bottomrule
\end{tabularx}
\end{table}
